\documentclass[../../../../main.tex]{subfiles}

\begin{document}

\section*{第1問}

\begin{proof}[解]
$C_1,L$，$C_2$と$L$の原点以外の交点$A(\alpha,k\alpha)$，$B(\beta,k\beta)$とおく。$\alpha,\beta$は各々，$x$の2次方程式$ax^2+bx=kx$，$px^2+qx=kx$の解のうち$x\ne0$のもので，
\[
\alpha=\frac{k-b}{a}, \qquad \beta=\frac{k-q}{p} \quad\cdots\text{①}
\]
である。
\[
S_1=\left|\int_\alpha^0(ax^2+bx-kx)dx\right|=\frac{|a|}{6}|\alpha|^3
\]
\[
S_2=\frac{|p|}{6}|\beta|^3
\]
だから，①から
\[
\frac{S_1}{S_2}=\frac{|a||\alpha|^3}{|p||\beta|^3}=\frac{|a|\cdot\frac{|k-b|^3}{|a|^3}}{|p|\cdot\frac{|k-q|^3}{|p|^3}}=\frac{|k-b|^3/a^2}{|k-q|^3/p^2} \quad\cdots\text{②}
\]
である。②が$k$によらない，すなわち$\dfrac{|k-b|^3}{a^2}=A\dfrac{|k-q|^3}{p^2}$が，$k$の恒等式になるような$A$の存在条件をかんがえれば良い。両辺0以上だから2乗して良く，
\[
p^4(k-b)^6=A^2a^4(k-q)^6 \quad\cdots\text{③}
\]
係数比較して，
\[
6\text{次}\cdots\ p^4=A^2a^4 \quad\cdots\text{④}
\]
\[
5\text{次}\cdots\ -6bp^4=A^2a^4(-6q) \quad\cdots\text{⑤}
\]
$A\ne0,p\ne0,a\ne0$だから，④⑤から，$b=q\ \cdots$⑥及び$p=\pm\sqrt Aa\ \cdots$⑦が必要。逆にこの時②において，
\[
\frac{|k-b|^3}{a^2}:\frac{|k-q|^3}{p^2}=\frac{|k-b|^3}{a^2}:\frac{|k-b|^3}{A\cdot a^2}=1:\frac1A \quad(\because k\ne b)
\]
で，たしかに$S_1:S_2$は一定で十分。①をみたす$A$は必ず存在するから，求める条件は⑥で
\[
b=q
\]
\end{proof}

\newpage

\section*{第2問}

\begin{proof}[解]
\textbf{(1)}
\[
\frac{\sin(2n+1)x}{\sin x}=\frac{x}{\sin x}\cdot\frac{\sin(2n+1)x}{(2n+1)x}\cdot(2n+1)\xrightarrow{x\to0}2n+1
\]

\textbf{(2)} $\displaystyle\int_0^{\frac{\pi}{2}}\frac{\sin(2n+1)x}{\sin x}dx=\frac{\pi}{2}\ \cdots$①を帰納法で示す。$n=1$の時，
\[
\int_0^{\frac{\pi}{2}}\frac{\sin3x}{\sin x}dx=\int_0^{\frac{\pi}{2}}(3-4\sin^2x)dx=\left[3x\right]_0^{\frac{\pi}{2}}-4\cdot\frac12\left[x-\frac12\sin2x\right]_0^{\frac{\pi}{2}}
\]
\[
=\frac32\pi-2\left(\frac{\pi}{2}\right)=\frac12\pi
\]
で成立。以下$n=k\in\mathbb{N}$での成立を仮定し，$n=k+1$での成立を示す。
\[
\int_0^{\frac{\pi}{2}}\frac{\sin(2k+3)x}{\sin x}dx=\int_0^{\frac{\pi}{2}}\frac{\sin(2k+1)x\cos2x+\sin2x\cos(2k+1)x}{\sin x}dx
\]
\[
=\int_0^{\frac{\pi}{2}}\frac{\sin(2k+1)x}{\sin x}(1-2\sin^2x)dx+\int_0^{\frac{\pi}{2}}2\cos x\cos(2k+1)x\,dx \quad\cdots\text{②}
\]
であり，
\[
\int_0^{\frac{\pi}{2}}\frac{\sin(2k+1)x}{\sin x}dx=\frac{\pi}{2} \quad(\because\text{仮定})
\]
\[
\int_0^{\frac{\pi}{2}}\sin x\sin(2k+1)x\,dx=0 \quad(\because k\ge1)
\]
\[
\int_0^{\frac{\pi}{2}}\cos x\cos(2k+1)x\,dx=0 \quad(\because k\ge1)
\]
だから，②に代入して
\[
\int_0^{\frac{\pi}{2}}\frac{\sin(2k+3)x}{\sin x}dx=\frac{\pi}{2}-0+0=\frac{\pi}{2}
\]
より，$n=k+1$でも①は成立し，よって示された。
\end{proof}

\newpage

\section*{第3問}

\begin{proof}[解]
$f(x)=x^4-2ax^2$とおく。$x=t$での接線は$\ell_t: y=(4t^3-4at)x-3t^4+2at^2$であるから，$C$と$\ell_t$の交点の$x$座標は（$P$のぞく）
\[
x^4-2ax^2=(4t^3-4at)x-3t^4+2at^2
\]
\[
x^4-2ax^2-4t(t^2-a)x+3t^4-2at^2=0
\]
\[
(x-t)^2\{x^2+2tx+(3t^2-2a)\}=0
\]
のうち$x\ne t$のもの，すなわち
\[
x^2+2tx+3t^2-2a=0 \quad(\text{左辺}g(x)\text{とする}) \quad\cdots\text{①}
\]
の解である。この判別式$D$として，$-\sqrt a\le t\le\sqrt a$から
\[
D/4=t^2-(3t^2-2a)=2a-t^2\ge0
\]
となり，①は異2実解（重解含む）を持つ。したがって，①の2解が$\alpha,\beta$で，
\[
\alpha+\beta=-2t, \qquad \alpha\beta=3t^2-2a
\]

\textbf{(2)} $\alpha\le\beta$だから，題意の条件は（重なる場合をのぞいて）
\[
\alpha<t<\beta \quad\cdots\text{②}
\]
である。すなわち，①の解が$x<t$に1つ，$t<x$に1つある時で，この条件は
\[
g(t)<0 \iff 6t^2-2a<0 \iff -\sqrt{\frac{a}{3}}<t<\sqrt{\frac{a}{3}}
\]
である。

\textbf{(3)} $\ell$の傾きは$4t^3-4at$だから，位置関係は右図のようになる。したがってピタゴラスの定理から
\[
L^2=(\beta-\alpha)^2+(\beta-\alpha)^2(4t^3-4at)^2
\]
\[
=\{1+16t^2(t^2-a)^2\}\{(\alpha+\beta)^2-4\alpha\beta\}
\]
\[
=\{1+16t^2(t^2-a)^2\}\{(-2t)^2-4(3t^2-2a)\} \quad(\because(1))
\]
\[
=\{1+16t^2(t^2-a)^2\}(8a-8t^2)
\]

\textbf{(4)} $L^2\equiv h(p)\ (p=t^2)$とすると，$-\sqrt a\le t\le\sqrt a$から$0\le p\le a\ \cdots$③である（(3)より）。
\[
h(p)=\{1+16p(p-a)^2\}(8a-8p)=-8(p-a)\{1+16p(p-a)^2\}
\]
だから，
\[
-\frac18h'(p)=1+16p(p-a)^2+(p-a)\cdot16(3p^2-4ap+a^2)
\]
\[
=16(p-a)\{p(p-a)+3p^2-4ap+a^2\}+1
\]
\[
=16(p-a)^2(4p-a)+1
\]
$a=\dfrac{7}{12}$の時，$p=\dfrac1{12}$を代入すると，$p-a=-\dfrac12,\ 4p-a=-\dfrac14$から
\[
16\left(-\frac12\right)^2\left(-\frac14\right)+1=16\cdot\frac14\cdot\left(-\frac14\right)+1=-1+1=0
\]
となり，$h'(1/12)=0$である。$0\le p\le a=7/12$での$h'(p)$の符号を調べると，下表のようになり，
\begin{center}
\begin{tabular}{c|ccccc}
$p$ & $0$ & & $1/12$ & & $7/12$ \\
\hline
$h'$ & & $+$ & $0$ & $-$ & \\
\hline
$h$ & & $\nearrow$ & & $\searrow$ &
\end{tabular}
\end{center}
したがって，$h(p)$は$p=\dfrac1{12}$でmaxで，
\[
\max L^2=-8\left(\frac1{12}-\frac{7}{12}\right)\left\{1+16\cdot\frac1{12}\left(\frac1{12}-\frac{7}{12}\right)^2\right\}
\]
\[
=4\left(1+\frac43\cdot\frac14\right)=\frac{16}{3}
\]
これから
\[
\max L=\sqrt{\frac{16}{3}}=\frac{4\sqrt3}{3}
\]
\end{proof}

\newpage

\section*{第4問}

\begin{proof}[解]
$n$次多項式を$P_n(x)$とおく。題意（$\diamondsuit$とおく）を帰納的に示す。$n=1$の時，$P_1(x)=ax+b\ (a\ne0)$とおけて，$P_1(0),P_1(1)\in\mathbb{Z}$から
\[
a+b\in\mathbb{Z}, \quad b\in\mathbb{Z} \quad\therefore\ a,b\in\mathbb{Z}
\]
となり，たしかに任意のセイスウ（整数）$k$に対し$P_1(k)\in\mathbb{Z}$となるから$\diamondsuit$は成立。以下$n\in\mathbb{N}$での$\diamondsuit$の成立を仮定し，$n=n+1$での成立を示す。そこで$P_{n+1}(x)$について，$P_{n+1}(0),\cdots,P_{n+1}(n+1)$が全て整数だとする。$Q(x)=P_{n+1}(x+1)-P_{n+1}(x)$とおくと，$Q(x)$は$n$次以下の多項式で，$Q(0),Q(1),\cdots,Q(n)$が全て整数なので，帰納法の仮定から，任意の$k\in\mathbb{Z}$に対して
\[
Q(k)\in\mathbb{Z}
\]
，すなわち
\[
P_{n+1}(k+1)=P_{n+1}(k)+Q(k)
\]
から，任意の$k$に対して$P_{n+1}(k)\in\mathbb{Z}$となり（$\because$仮定），$n=n+1$でも$\diamondsuit$は成立。以上から示された。
\end{proof}

\newpage

\section*{第5問}

\begin{proof}[解]
$x_k=1,2,\cdots,6$だから（$k=1,2,3,4$），$P,O,Q$の3点は常に異なるので，
\[
\angle POQ\text{が鋭角} \iff \overrightarrow{OP}\cdot\overrightarrow{OQ}>0 \iff x_1x_3>x_2x_4 \quad\cdots\text{①}
\]
である。ここで，さいころの出る目6通りのうち，
\[
\begin{cases}
x_1x_3>x_2x_4 & \cdots A\text{通り} \\
x_1x_3=x_2x_4 & \cdots B\text{通り} \\
x_1x_3<x_2x_4 & \cdots C\text{通り}
\end{cases}
\]
とおくと，対称性から$A=C\ \cdots$②であり，上で全ての場合が尽くされ，かつ排反だから
\[
A+B+C=6^4 \quad\cdots\text{③}
\]
も成立する。ここで，$B$を数える。$p,q\ (p,q=1,2,\cdots,6)$に対する積$pq$は，以下の表になる。
\begin{center}
\begin{tabular}{c|cccccc}
$p\backslash q$ & 1 & 2 & 3 & 4 & 5 & 6 \\
\hline
1 & 1 & 2 & 3 & 4 & 5 & 6 \\
2 & 2 & 4 & 6 & 8 & 10 & 12 \\
3 & 3 & 6 & 9 & 12 & 15 & 18 \\
4 & 4 & 8 & 12 & 16 & 20 & 24 \\
5 & 5 & 10 & 15 & 20 & 25 & 30 \\
6 & 6 & 12 & 18 & 24 & 30 & 36
\end{tabular}
\end{center}
積$v$を実現する$(p,q)$の組の個数を$n_v$とすると，$x_1x_3=x_2x_4=v$となる場合の数は，$(x_1,x_3)$，$(x_2,x_4)$ともに積$v$を実現する組であればよいから，$n_v^2$通り。表から，$n_v=1$となる積（1,9,16,25,36の5個），$n_v=2$となる積（2,3,5,8,10,15,18,20,24,30の10個），$n_v=3$となる積（4の1個），$n_v=4$となる積（6,12の2個）に分けて，
\[
B=\sum_vn_v^2=5\cdot1^2+10\cdot2^2+1\cdot3^2+2\cdot4^2=5+40+9+32=86
\]
だから，②③から
\[
A=\frac{6^4-86}{2}
\]
となり，①から，もとめる確率は
\[
\frac{\frac{6^4-86}{2}}{6^4}=\frac{3\cdot6^3-43}{6^4}=\frac{605}{1296}
\]
\end{proof}

\end{document}
